% $Author: stef $
% $Date: 2008-04-04 17:14:31 +0200 (Fri, 04 Apr 2008) $
% $Revision: 318 $
%=================================================================
\ifx\wholebook\relax\else
% --------------------------------------------
% Lulu:
    \documentclass[a4paper,10pt,twoside]{book}
    \usepackage[
        papersize={6in,9in},
        hmargin={.75in,.75in},
        vmargin={.75in,1in},
        ignoreheadfoot
    ]{geometry}
    \input{../common.tex}
    \pagestyle{headings}
    \setboolean{lulu}{true}
% --------------------------------------------
% A4:
%   \documentclass[a4paper,11pt,twoside]{book}
%   \input{../common.tex}
%   \usepackage{a4wide}
% --------------------------------------------
    \graphicspath{{figures/} {../figures/}}
    \begin{document}
%   \renewcommand{\nnbb}[2]{} % Disable editorial comments
    \sloppy
\fi

\chapter{Composing Messages}\label{cha:composingMessages}

%%%%%%%%%%%%%%%%%%%%%%%%%%%%%%%%%%%%%%%%%%%%%%%%%%%%%%%%%%%%%%%%%%%%%%%%
% \begin{figure}[h]
% 	\centerline{\includegraphics{}}
% 	\caption{Ready-to-use files for the Macintosh.  
% 	\label{fig:macfiles}}
% \end{figure}

% %%%%%%%%%%%%%%%%%%%%%%%%%%%%%%%%%%%%%%%%%%%%%%%%%%%%%%%%%%%%%%%%%%%%%%%%
% \begin{exofigwithsize}[0.5]{\includegraphics[width=3cm]{turtleMSmallStairs}}{A Staircase}\label{xp:letterA}
% You are not limited in your robot drawings to squares. You can create a wide range of geometrical figures.
% For example, here is a drawing of a small staircase. Write a script to reproduce this drawing. 
% \end{exofigwithsize}
% 
% \begin{exonofigtitle}{Moving Clock Hands}
% Experiment with different angle values for each of the two robots; that is,change the angle values for the two turn 
% methods. Then, compare the effect of the method \ct{turnLeft: 60} (for pica) and \ct{turnRight: 300} (for daly). 
% You can see that turning left 60 degrees yields the same result as turning right 300 degrees. This is so because 
% the sum of the two values is 360 degrees, that is, a full circle. 
% \end{exonofigtitle}
% 
% %%%%%%%%%%%%%%%%%%%%%%%%%%%%%%%%%%%%%%%%%%%%%%%%%%%%%%%%%%%%%%%%%%%%%%%%
% \begin{scriptfigwithsize}[0.4]{\includegraphics[width=5cm]{turtleMLetterA}}{The letter A}\label{scr:letterA}
% 	| pica | 
% 	pica := Bot new. 
% 	pica north. 
% 	pica go: 100. 
% 	pica east. 
% 	pica go: 100. 
% 	pica south. 
% 	pica go: 100. 
% 	pica north. 
% 	pica go: 50. 
% 	pica west. 
% 	pica go: 100
% \end{scriptfigwithsize}
% 
% \begin{script}[myster1]{What does pica do? (Problem 1)}
% 	| pica | 
% 	pica := Bot new. 
% 	pica go: 100. 
% 	pica turnLeft: 45. 
% 	pica go: 50. 
% 	pica turnLeft: 45. 
% 	pica go: 100 
% \end{script}

As with any language, Smalltalk follows certain rules in executing the messages sent to 
objects. I have not yet presented these rules to you, and you may have wondered just what 
those rules are when you were experimenting with the previous scripts. Your patience will be 
now rewarded. This chapter explains how to read and write correctly formulated messages. 
This chapter may appear a bit more difficult or abstract than the previous ones. However, 
I have done my best to present clearly the simple rules that govern the writing of messages. 
Understanding such details might not be as much fun as playing with robots, but it is the 
price that must be paid if you are to be able to write more advanced programs. The good news 
is that Smalltalk is not a complex language: there are only five rules that you need to under- 
stand. If you are still hesitant, you may also skip this chapter on a first reading and return to it 
when you have questions about the structure of your programs. 

As described in Chapter 2, a message is composed of the \emph{message selector} and the optional 
\emph{message arguments}. A message is sent to a \emph{message receiver}. The combination of a message and 
its receiver is called a \emph{message send}. 

When you write a complex expression such as pica \ct{go: 100 + 20}, it contains two message 
sends, using the message selectors \ct{go:} and \ct{+}, and you have to know the order in which the 
messages are executed if you are to understand what the result of the entire expression will be. 
In Smalltalk, the order in which messages are executed is determined by the type of message 
send. There are three types of messages: \emph{unary}, \emph{binary}, and \emph{keyword-based}. Unary messages 
are always sent first, followed by binary messages, and finally keyword-based messages. Any 
messages enclosed in parentheses are executed prior to any other messages. This means that 
you can change the order in which message sends are executed through the use of parentheses. 
These rules go a long way toward making Smalltalk code easy to read. And you will soon 
discover that most of the time, you do not even have to think about them. However, you have 
to know them, because occasions will arise that will require your knowledge of them. 

All of the examples presented in this chapter consist of executable code, as shown in the 
text. So do not hesitate to try them out and see how they function. I will begin by showing you 
how to identify the different types of messages, and then I will present some examples of each 
type. Finally, I will present the rules for message composition. 



\section{The Three Types of Messages }

Smalltalk defines a small number of simple rules to determine the order in which message 
sends are executed. These rules, which I will present in detail later, are based on the distinction 
among three different types of messages: 

\begin{itemize}
\item  \emph{Unary messages} are messages with no arguments. They are sent to an object (the 
message receiver) without any other information. For example, in the expression \ct{pica color}, 
the message color is a unary message. It does not send any additional information; 
that is, there is no argument. These messages are called “unary,” from the Latin 
\emph{unum}, meaning “one,” because a message send with a unary message involves only 
\emph{one object}, the message receiver. 

\item  \emph{Binary messages} are messages that involve two objects: the message receiver and the 
message selector’s sole argument. (The word “binary” is related to the Latin \emph{bis}, mean- 
ing “twice.”) Binary messages are mainly related to mathematical expressions. For 
example, in the message send \ct{10 + 20}, the message consists of the message selector + 
together with the single argument 20. The message \ct{+ 20} is sent to the object 10, which 
is the message receiver. 

\item  \emph{Keyword-based messages} are messages whose message selector contains a keyword 
with at least one colon character \ct{:} in its name and that has one or more arguments. For 
example, in the message send \ct{pica go: 100}, the keyword \ct{go:} contains the colon character: and 
there is one argument, \ct{100}, and therefore there are two objects involved in 
the message send: the message receiver \ct{pica} and the argument \ct{100}. 
\end{itemize}


\important{Don’t be confused by the nomenclature for unary and binary messages. The idea of “one” in the 
word \emph{unary} and the idea of “two” in the word \emph{binary} refer to the number of objects involved in a message 
send, not the number of arguments. Thus a unary message has no arguments, and so a unary message 
send involves one object, namely,the message receiver. A binary message has a single argument,and 
therefore a binary message send involves two objects: the argument of the message selector and the 
message receiver.}	

\section{Identifying Messages}

In order to understand the structure of an expression, the first thing you need to do is to identify 
the messages and their receivers. To do this, I suggest that you use a graphical notation as shown 
in Figure~\ref{fig:uKeyUnOne}. In all the figures of this chapter, the message receivers are underlined, each 
message is surrounded by an \emph{ellipse}, and the messages that make up the expression are numbered in 
the order in which they will be executed. When there is more than one ellipse, the first ellipse to 
be executed is drawn with a dashed line so that you can see at once where to begin. 

Figure~\ref{fig:uKeyUnOne}  shows that the expression \ct{pica color: Color yellow} contains within it the 
expression \ct{Color yellow}. Therefore, there are two ellipses, one for the entire expression \ct{pica color: Color yellow}, and one for the subexpression \ct{Color yellow}. You will learn a bit later that 
the expression \ct{Color yellow} is executed first, so its ellipse is a broken line and is numbered 1. 

Note that each message has a receiver. The robot pica receives the message \ct{color: ...}, and 
Color receives the message yellow. Therefore, these two message receivers are underlined. (The 
three dots in the message \ct{color: ...} indicate the argument of the message selector \ct{color:}, 
which will be the result of the message send \ct{Color yellow}.) 

 \begin{figure}[h]
 	\centerline{\includegraphics[width=5cm]{uKeyUnOne}}
 	\caption{The expression \ct{pica color: Color yellow} contains the subexpression \ct{Color yellow}. 
	The message \ct{yellow} is sent to \ct{Color}, and then the message \ct{color: ...} is sent to \ct{pica}.  
 	\label{fig:uKeyUnOne}}
 \end{figure}

As I have mentioned, every message is sent to an object called the message receiver. A 
receiver does not have to be a robot. It can be just about anything, from a number to a window. A message receiver can appear explicitly as the first element of an expression, such as 
\ct{pica} in the expression \ct{pica go: 100} or \ct{Color} in the expression \ct{Color new}. However, a receiver 
can also be the result of other message sends. For example, in the message \ct{Bot new go: 100}, 
the receiver of the message \ct{go: 100} is the resulting object returned by the message send Bot 
new. Nevertheless, in every case, a message is sent to some object, and that object is called the 
message receiver. 

\important{In a message send, a message is always sent to an object, called the message receiver, 
which may be named explicitly or may be the result of other message sends.}

As I have mentioned, every message is sent to an object called the message receiver. A 
receiver does not have to be a robot. It can be just about anything, from a number to a win- 
dow. A message receiver can appear explicitly as the first element of an expression, such as 
\ct{pica} in the expression \ct{pica go: 100} or \ct{Color} in the expression \ct{Color new}. However, a receiver 
can also be the result of other message sends. For example, in the message \ct{Bot new go: 100}, 
the receiver of the message \ct{go: 100} is the resulting object returned by the message send Bot 
new. Nevertheless, in every case, a message is sent to some object, and that object is called the 
message receiver. 


Table~\ref{tab1101} shows some message sends, and I suggest that in each case, you identify the 
type of message and then draw the graphical representation of the expression.

 \begin{table}[h]
 \caption{\emph{Some Examples of Message Sends, Simple and Compound.}\label{tab1101}}
 \begin{center}
 {\small \begin{tabular}{p{23mm}p{28mm}p{50mm}}
 \hline
 \textbf{Expression} & \textbf{Type(s)} & \textbf{Action}\\ \hline
 \textsf{pica go: 100} & keyword-based & The receiving robot moves forward 100 pixels. \\
 \textsf{100 + 20} & binary & The number 100 receives the message with the number 20 as argument.\\
 \textsf{pica est} & unary & The receiving robot \textsf{pica} points to the \textsf{east}. \\
 \textsf{pica color: Color yellow} & keyword-based, unary & The receiving robot \textsf{pica} changes its color to the color yellow.\\
 \textsf{pica go: 100 + 20} & keyword-based, binary & The receiving robot moves forward 120 pixels.\\
 \textsf{Bot new go: 100} & unary, keyword-based & The message \textsf{new} is sent to the Bot class, 
 which returns a new robot to which the message \textsf{go: 100} is sent, causing the robot to move forward 100 pixels. \\
 \hline
 \end{tabular}}
 \end{center}
 \end{table}



From Table~\ref{tab1101} you should observe the following points: 
\begin{itemize} 
	
	\item Some of the messages have arguments, while others do not. The message east, being 
unary, does not have an argument, while \ct{go: 100} and \ct{+ 20} each have one argument, 
the number 100 and the number 20, respectively. 
\item  Different messages are sent to different objects. In the expression pica east, the message
\ct{east} is sent to a robot, and in \ct{100 + 20}, the message \ct{+ 20} is sent to the number \ct{100}. 

\item  There are simple messages and compound messages. For example, \ct{Color yellow} and 
\ct{100 + 20} are simple: One message is sent to one object. On the other hand, the expression 
\ct{pica go: 100 + 20} contains two messages: \ct{+ 20} is sent to \ct{100}, and then \ct{go: ...} is sent to 
pica, where ... represents the result of the execution of the message send \ct{100 + 20}. 
\item  A message receiver can be the result of an expression that returns an object. In the 
expression \ct{Bot new go: 100}, the message \ct{go: 100} is sent to the object (a robot) that 
results from evaluation of the expression \ct{Bot new}. 
\end{itemize}

\section{The Three Types of Messages in Detail}
Now that you are able to identify message receivers and the three types of messages, let us 
look at the different types of messages in detail. 


\subsection{Unary Messages}

Among the uses of unary messages are obtaining a value from an object, such as the size of a 
robot’s pen (\ct{pica penSize}); obtaining an object from a class (\ct{Color yellow}); and instructing 
the receiver to perform an action (\ct{pica beInvisible}). Recall that a unary message does not 
take an argument: it is sent to the receiver without any other information. Thus the one object 
involved in a unary message send is the message receiver. Unary message sends are of the 
form of \ct{receiver messageName}. Script~\ref{scr:unaryex} presents some examples of unary messages, shown 
in boldface type. 

\begin{script}[unaryex]{Examples of unary messages}
| pica | 
pica := Bot new. 
pica !\textbf{color}!. 
pica !\textbf{penSize}!. 
pica !\textbf{east}!. 
Color !\textbf{yellow}!. 
125 !\textbf{factorial}!
\end{script}





\important{Unary messages are messages that do not take an argument. A unary message send has 
the form \textsf{receiver messageName}.}

\subsection{Binary Messages}

Binary messages are messages that involve two objects: the receiver and a single argument. All 
binary message selectors are composed of one or two characters from the following list: 
\ct{+}, 
\ct{*},
\ct{/}, 
\ct{|},
\ct{&},
\ct{=},
\ct{>},
\ct{<},
\textsf{$\sim$},
\ct{@}. Therefore \ct{+}, \ct{=}, and \ct{*} are message selectors, but so is \ct{=>}, which is composed of two symbols. 

Table~\ref{tab1102} shows some examples of binary message sends and their meaning. At this 
point, I would rather not go into the details of these examples, so don’t worry if you are not 
sure about exactly what each of these message selectors does. But try executing the expressions and others like them. 


\begin{table}[h]
 \caption{\emph{Examples of Binary Messages with Numbers.}\label{tab1102}}
 \begin{center}
 {\small \begin{tabular}{p{23mm}p{28mm}p{50mm}}
 \hline
 \textbf{Expression} & \textbf{Returned Value} & \textbf{Action}\\ \hline
\textsf{1 + 2.5}& \textsf{3.5}& Addition of two numbers \\ 
\textsf{3.4 * 5}& \textsf{17.0}& Multiplication of two numbers \\ 
\textsf{8 / 2}& \textsf{4}& Division of two numbers \\ 
\textsf{10 - 8.3}& \textsf{1.7}& Subtraction of two numbers \\ 
\textsf{12 = 11}& \textsf{false}& Testing for equality between two numbers \\ 
\textsf{12 ~= 11}& \textsf{true}& Testing for inequality between two numbers \\ 
\textsf{12 > 9}& \textsf{true}& Is the receiver greater than the argument? \\ 
\textsf{12 >= 10}& \textsf{true}& Is the receiver greater than or equal to the argument? \\ 
\textsf{12 < 10}& \textsf{false}& Is the receiver less than the argument? \\ 
\textsf{100@10}& \textsf{100@10}& Create a point with coordinates (100, 10)\\ 
 \hline
 \end{tabular}}
 \end{center}
 \end{table}



\important{Binary message sends involve two objects:the receiver and a single argument.The mes- 
sage selector of a binary argument is composed of one or two characters from the following list: \ct{+}, 
\ct{*}, \ct{/}, \ct{|}, \ct{&}, \ct{=}, \ct{>}, \ct{<}, \textsf{$\sim$}, \ct{@}. Binary message sends have the form \ct{receiver messageName argument}.}



\subsection{Keyword-Based Messages}

Keyword-based messages are messages that take at least one argument and that contain at least 
one colon character \ct{:}. Note that the colon is part of the message selector. Therefore, \ct{go:}, not \ct{go},  is the name of a keyword-based message. Script~\ref{scr:unarykey} shows some examples of keyword-based 
messages, shown in boldface type. 



\begin{script}[unarykey]{Examples of keyword-based messages}
	| pica | 
	pica := Bot new. 
	pica !\textbf{go:}! 100. 
	pica !\textbf{penSize:}! 5. 
	pica !\textbf{color:}! Color yellow. 
	pica !\textbf{turn:}! 90
\end{script}


I have said that a keyword-based message has at least one argument, but we have not yet 
seen an example of such a message with multiple arguments. Let us look at an example now: 
The message send \ct{aNumber between: lowerBound and: upperBoundchecks} whether the number 
\ct{aNumber} is in the interval represented by the two numbers \ct{lowerBound} and \ct{upperBound}. 
This message needs two arguments, namely, the two bounds of the interval. An example is 
shown in Table 11-3. Note that the message selector is actually \ct{between:and:}. It is composed 
of the two words \ct{between:} and \ct{and:}. 


\begin{table}[h]
 \caption{\emph{Keyword-Based Messages That Take More Than One Argument.}\label{tab1103}}
 \begin{center}
 {\small \begin{tabular}{p{25mm}p{18mm}p{20mm}p{40mm}}
 \hline
 \textbf{Expression} & \textbf{Arguments} &\textbf{Returned Value} & \textbf{Action}\\ \hline
5 between: 2 and: 10 &2, 10&true &Is 5 between 2 and 10? \\
Color r:0g:1b:0 &0, 1, 0&a green color object&creates a color with the given values of red, green, and blue. \\ \hline
 \end{tabular}}
 \end{center}
 \end{table}


\important{Keyword-based messages have at least one argument, and their message selector contains 
at least one colon character \ct{:}. A keyword-based message send that takes two arguments is of the form 
\ct{receiver messageNameWordOne: argumentOne messageNameWordTwo: argumentTwo}.}

\section{Order of Execution}
You have seen that there are three kinds of messages: unary, binary, and keyword-based. Now 
I will tell you, as I promised, how to determine the order in which messages are executed. The 
order of message execution is determined by the type of message, as described by the follow- 
ing three rules: 

\begin{description}
\item Rule 1: Unary message sends are executed first, then binary message sends, and finally 
keyword-based ones. 
\item Rule 2: As with mathematical expressions, the priority of message execution can be over- 
ridden by parentheses: message sends in parentheses are executed before any other types of message sends. 
\item Rule 3: Message sends of the same type are executed from left to right. 
\end{description}

These rules may seem complex, but they are quite natural, and once you get used to them, 
you will not have to think too much about them most of the time. In particular, the third rule 
simply states that messages of the same type are executed in the order in which they are read. 

If you are ever in doubt and want to be sure that your messages are executed the way you 
want, you can always add extra parentheses, as shown in Figure~\ref{fig:uKeyUn}. 

\begin{figure}[h]
	\centerline{\includegraphics[width=\textwidth]{uKeyUn}}
	\caption{Unary message sends are executed first, so \ct{Color yellow} is executed first. This execution returns a color object,which is passed as the argument of the message \ct{color:} ... that is 
	sent to \ct{pica}. 
	\label{fig:uKeyUn}}
\end{figure}


In the figure, the expression \ct{pica color: Color yellow} is analyzed. The message selector \ct{yellow} is a unary message, while the message selector \ct{color:} is a keyword-based one. Therefore, the expression \ct{Color yellow} is executed first. If you are unsure about the order of execution, then you can put
parentheses around \ct{Color yellow} to make sure that it will be executed first. This won’t change 
the natural order of execution. That is, \ct{pica color: Color yellow} and \ct{pica color: (Color yellow)} have precisely the same effect. The rest of this section illustrates each of these points. 


\section{Rule 1: Unary > Binary > Keywords}

Unary message sends are executed first, then binary message sends, and finally keyword- 
based message sends. In programmer jargon we also say that unary messages have precedence 
over binary messages, and binary messages have precedence over keyword-based messages. 
\important{Rule 1:Unary message sends are executed before binary message sends,which are 
executed before keyword-based message sends.}

\subsection{Example 1}
In the message send \ct{pica color: Color yellow}, there is one unary message, \ct{yellow}, sent 
to the class \ct{Color}, and there is one keyword-based message, \ct{color:} ..., which is sent to the 
robot \ct{pica}. Unary message sends are executed first, so the first message send to be executed 
is \ct{Color yellow}. This execution returns a color object, represented as \ct{aColor} (since we need a 
name to refer to it), which is passed to pica as the argument of the message \ct{color: aColor}. 
Figure \ref{fig:uKeyUn} shows graphically the order in which the messages are executed. 

As an aid to your understanding, I would like to propose a textual way of representing a 
compound message send in step-by-step execution. In Step-by-step~\ref{scr:unarykey2}, the message send 
to be executed step by step is \ct{pica color: Color yellow}. The first line shows the complete 
message send in boldface type.

\begin{script}[unarykey2]{Step-by-step decomposition of the execution of pica color: Color yellow}
pica color: Color yellow 
(1)           Color yellow        "unary" 
-returns> aColor  
(2) pica color: aColor "keyword-based" 
\end{script}

The lines of code represent the execution steps in numbered order in which they will occur. 
Thus \ct{Color yellow} is the first expression to be executed. Note that the expressions have been 
indented to line up with their counterparts in the message send at the top. 

When the execution of a message send returns a result that is used in the following execution, the line following the executed expression shows “-returns>” followed by the result. Here 
the expression \ct{Color yellow} returns a color object that I have called \ct{aColor} so that it can be 
referred to in the sequel. The second expression to be executed is \ct{pica color: aColor}, where 
as I just explained, \ct{aColor} is the result obtained from the previous execution step. To stress 
this point, the returned value or object is displayed in italic type. For further clarification, the 
kind of message that is currently being executed is displayed as a comment inside quotation 
marks. For example, \ct{Color yellow} is shown to be a unary message. 


\subsection{Example 2}
The message send \ct{pica go: 100 + 20}, contains the binary message selector + and the keyword-based message selector \ct{go:}. Binary messages are executed prior to keyword-based 
messages, so \ct{100 + 20} is executed first: The message \ct{+ 20} is sent to the object \ct{100}, which 
returns the number 120. Then the message \ct{pica go: 120} is executed with \ct{120} as argument. 
Step-by-step script \ref{scr:uKeyBin} shows how the expression is executed (Figure~\ref{fig:uKeyBin}). 

\begin{figure}[h]
	\centerline{\includegraphics[width=5cm]{uKeyBin}}
	\caption{Decomposition of the expresssion \ct{pica go: 100 + 20}. 
	\label{fig:uKeyBin}}
\end{figure}

\begin{script}[uKeyBin]{Step-by-step decomposition of the expresssion pica go: 100 + 20}
pica go: 100 + 20 
(1)          100 + 20      "binary" 
-returns> 120 
(2) pica go: 120           "keyword-based" 
\end{script}





\subsection{Example 3}
The message \ct{pica penSize: pica penSize + 2} contains the unary message \ct{penSize}, the 
binary message with selector \ct{+}, and the keyword-based message with selector \ct{penSize:}. 
Step-by-step script \ref{scr:uKeyUnBin} illustrates the decomposition of message execution. The unary message send 
\ct{pica penSize} is executed first (step 1). This message returns a number, which we are calling 
\ct{aNumber}, representing the current size of the receiver’s pen. Then the binary message send 
\ct{aNumber + 2} is executed (step 2). The number \ct{aNumber} is the receiver of the message \ct{+ 2}, 
which in turn returns another number, the sum, which here is called \ct{anotherNumber}. Finally, 
the keyword-based message \ct{penSize: anotherNumber} is sent to \ct{pica}, who sets his pen size to 
\ct{anotherNumber} (see Figure~\ref{fig:uKeyUnBin}). 

\begin{figure}[h]
	\centerline{\includegraphics[width=5cm]{uKeyUnBin}}
	\caption{Decomposition of the expresssion \ct{pica penSize: pica penSize + 2}. 
	\label{fig:uKeyUnBin}}
\end{figure}

\begin{script}[uKeyUnBin]{Step-by-step decomposition of the expresssion pica go: 100 + 20}
pica penSize: pica penSize + 2 
(1)               pica penSize        "unary" 
-returns> aNumber  
(2)               aNumber+ 2         "binary" 
-returns> anotherNumber 
(3) pica penSize: anotherNumber "keyword-based" 
\end{script}


Altogether, the entire compound expression increases the receiver’s pen size by two 
pixels. It does so by first asking \ct{pica} for his pen size (\ct{pica penSize}), increasing that number 
by two \ct{(aNumber+ 2)}, and then telling pica to change his pen size to the new number 
(pica penSize: anotherNumber). Note that \ct{penSize} and \ct{penSize:} are two different message 
selectors! The first is unary and asks the receiver for its pen size, and the second is keyword- 
based and tells the receiver to change its pen size to the value of its argument. 

\subsection{Example 4}
As an exercise, I will let you decompose the execution of the message \ct{Bot new go: 100 + 20}, 
which is composed of one unary, one keyword-based, and one binary message (see Figure~\ref{fig:uunKeyBin}). 

\begin{figure}[h]
	\centerline{\includegraphics[width=5cm]{uunKeyBin}}
	\caption{Decomposition of the expresssion \ct{pica  new go: 100 + 20}. 
	\label{fig:uunKeyBin}}
\end{figure}




\section{Rule 2: Parentheses First}

The default ordering of message execution may not be suitable for what you want to accomplish in 
an expression, and so you should be able to change it. For this purpose, Smalltalk 
offers parentheses \ct{(}and\ct{)}. Just as in mathematics, expressions in parentheses get the highest 
precedence, and they are executed before any others. 

Keep in mind that if you find the rules for order of execution a bit complex or if you sim- 
ply want to clarify the structure of an expression, use parentheses to ensure that the messages 
are executed in the order that you wish. Figure~\ref{fig:equi} shows some of the expressions that we 
have previously looked at together with their equivalents using parentheses. 

\begin{figure}[h]
	\centerline{\includegraphics[width=\linewidth]{uKeyUnBinPar}}
	\centerline{\includegraphics[width=\linewidth]{uunKeyBinPar}}
	\caption{Equivalent messages using parentheses.\label{fig:equi}}
\end{figure}


\important{Rule 2: As in mathematics, messages in parentheses are executed before any others.
They  have the highest priority.}

\subsection{Example 5}

The message \ct{(65 @ 325 extent: 134 @ 100) center} returns the center of a rectangle whose 
upper-left point has coordinates \ct{(65, 325)}, whose width is \ct{134} pixels, and whose height
 is \ct{100} pixels. Step-by-step~\ref{scr:thisone} shows how the message is decomposed and executed. First, the message 
within the parentheses is executed. It is a compound expression consisting of three message 
sends: two binary message sends, \ct{65 @ 325} and \ct{134 @ 100}, which are executed first and return 
points; and the keyword-based message \ct{extent:} ..., which is then sent to the point (65, 325) and 
which returns a rectangle with upper-left point and dimensions as described above. Finally, the 
unary message center is sent to this rectangle, and a point, the rectangle’s center, is returned. Try- 
ing to evaluate the message without parentheses leads to an error, because in that case, the unary 
message \ct{center} would have to be executed first, and it would be sent to the object \ct{100}, but a 
number object does not understand the message \ct{center}. 


\begin{script}[thisone]{Step-by-step decomposition of the expresssion \ct{pica go: 100 + 20}}
(65 @ 325 extent:  134 @ 100) center 
(1)  65@325                                      "binary" 
-returns> aPoint 
(2)                    134@100                   "binary" 
-returns> anotherPoint 
(3)  aPointextent: anotherPoint "keyword-based" 
-returns> aRectangle  
(4)  aRectanglecenter                           "unary" 
-returns> 132@375  
\end{script}



\section{Rule 3: From Left to Right}

Now that you know how messages are categorized according to priority of execution, the final 
question to be addressed is how messages with the same priority are executed. Rule 3 states 
that they are executed from left to right. This rule was used already, in Step-by-step~\ref{scr:thisone}, 
where the left-hand binary message \ct{@ 325} was executed before the right-hand message \ct{@ 100}. 

\important{Rule 3: Messages of the same type are executed in order from left to right.}

\subsection{Example 6}

In the expression \ct{Bot new east}, both message sends are unary messages, so the first one as 
you read from left to right, \ct{Bot new}, is executed first. It returns a newly created robot, called 
\ct{aBot} in Step-by-step~\ref{scr:noThatone} to which the second message, east, is sent. 
Figure~\ref{fig:ucompoUn} shows  the order of execution. 


\begin{script}[noThatone]{Step-by-step decomposition of the expresssion \ct{Bot new east}}
!\textbf{Bot new east}!     
  (1)  Bot new                 "unary"
      -returns>  aBot
  (2)         aBot east        "unary"
\end{script}


\begin{figure}[h]
	\centerline{\includegraphics[width=\linewidth]{ucompoUn}}
	\caption{Decomposition of the expression \ct{Bot new east}.\label{fig:ucompoUn}}
\end{figure}


\subsection{Example 7}

In the expression \ct{20 + 2 * 5}, there are only the two binary message selectors \ct{+} and \ct{*}. According 
to Rule 3, since \ct{+} is to the left of \ct{*}, it should be executed first. In normal mathematical 
notation as well as in many programming languages, multiplication would take precedence 
over addition regardless of the order in which the arithmetic operations appear. However, in 
Smalltalk there is \emph{no specific priority} for mathematical operations. The message descriptors \ct{+} 
and \ct{*} are just binary messages, and therefore they have equal status. The message selector \ct{*} 
does not have precedence over \ct{+}, and the leftmost message selector \ct{+} is sent first, and then 
\ct{*} is sent to the result, as shown in Step-by-step~\ref{scr:math} and Figure~\ref{fig:ucompoNoBracket}. 

\begin{script}[math]{Step-by-step decomposition of the expression \ct{20 + 2 * 5}}
!\textbf{20 + 2 * 5}!
  
(1)  20 + 2
  -returns> 22
(2)  22 * 5
   -returns> 110
\end{script}


\begin{figure}[h]
	\centerline{\includegraphics[width=\linewidth]{ucompoNoBracket}}
	\caption{Decomposition of the expression \ct{20 + 2 * 5}.\label{fig:ucompoNoBracket}}
\end{figure}


\important{There is no priority among binary messages. In the expression \ct{20 + 2 * 5}, the leftmost message 
\ct{+} is evaluated first, despite the fact that in normal mathematical notation, 
multiplication takes precedence over addition.}

You can see, then, as shown in Step-by-step~\ref{scr:math}, that the result of this expression is not 
30, which you would get if you did the multiplication first, but 110. This behavior is surprising 
at first, but it derives from the three simple rules for executing messages. This counter intuitive 
order of mathematical operations is the price that we have to pay for the simplicity of the 
Smalltalk model, which has only methods. If you want your expression to obey the normal 
priority of mathematical operations, then you should use parentheses, since when message 
sends are enclosed in parentheses, they are executed first. Hence the expression \ct{20 + (2 * 5)} 
returns the result 30, as shown in Step-by-step~\ref{scr:math2} and Figure~\ref{fig:ucompoNumberBracket}. 

\begin{script}[math2]{Step-by-step decomposition of the expression \ct{20 + (2 * 5)}}
!\textbf{20 + (2 * 5)}!
  
(1)  2 * 5
  -returns> 10
(2)  20 + 5
   -returns> 30
\end{script}


\begin{figure}[h]
	\centerline{\includegraphics[width=\linewidth]{ucompoNumberBracket}}
	\caption{Decomposition of the expression \ct{20 + (2 * 5)}.\label{fig:ucompoNumberBracket}}
\end{figure}


\important{Message sends surrounded by parentheses are executed first. Therefore,in the expression 
\ct{20 + (2 * 5)}, the message with \ct{*} is executed before the one with \ct{+}, which is the usual order of 
operations in mathematics.}

\important{
Note In Smalltalk,the mathematical message selectors such as \ct{+}and \ct{*} all have the same priority. The 
symbols \ct{+} and \ct{*} are simply message selectors for binary messages. Therefore, * does not have priority over 
\ct{+}. If you want to force one operation to take precedence over another, then you should use parentheses. 
The fact that Smalltalk does not follow mathematical precedence can be confusing at the beginning. Therefore, 
when you have multiple binary messages representing a mathematical expression, give yourself and 
anyone else reading your program a break and insert parentheses to express how the computation should 
be performed. When you have become more accustomed to the way that messages are executed, you will 
probably become more laissez-faire about parentheses.}

A consequence of rule 1 ---which provides for the order of execution of different types of 
messages, with unary messages being executed before binary messages, and binary messages 
before keyword-based messages---is that you very often do not have to use parentheses. That 
is, most of the time, you do not have to worry about order of execution. Table~\ref{tab1104} shows 
expressions written to be executed according to Smalltalk’s rules and equivalent expressions 
using parentheses, which would be necessary if the order of precedence did not exist. 


\begin{table}[h]
\caption{Some Expressions and Their Fully Parenthesized Equivalents.}
\label{tab1104}
\begin{center}
{\small \begin{tabular}{p{50mm}p{50mm}}
\hline
\textbf{Without Parentheses} & \textbf{Equivalent Expression with Parentheses} \\
\hline
\textsf{pica color: Color yellow} & \textsf{pica color: (Color yellow)} \\
\textsf{pica go: 100 + 20} & \textsf{pica go: (100 + 20)} \\
\textsf{pica penSize: pica penSize + 2} & \textsf{pica penSize: ((pica penSize) + 2)}\\
\textsf{2 factorial + 4} & \textsf{(2 factorial) + 4} \\
\hline
\end{tabular}}
\end{center}
\end{table}








\section{Summary}

\begin{itemize}
\item  A message is always sent to an object, called the message receiver, which may be the 
result of other messages. 
\item  Unary messages are messages that take no argument. A unary message send is of the 
form of \ct{receiver messageName}. 
\item  Binary messages are messages that involve two objects: the receiver of the message 
and a single argument. The message selector of a binary message consists of one or two 
characters from the following list: \textsf{$+$}, \textsf{$*$}, \textsf{$/$}, \textsf{$|$}, \textsf{\&}, \textsf{$=$}, \textsf{$>$}, \textsf{$<$}, \textsf{$\sim$}, \textsf{@}. Binary message sends are of 
the form \ct{receiver messageName argument}. 
\item  Keyword-based messages are messages that take one or more arguments and use a 
keyword with at least one colon character :. A keyword-based message send taking 
two arguments is of the form \ct{receiver messageNameWordOne: argumentOne messageNameWordTwo: argumentTwo}. 

\item Rule 1. Unary messages are executed first, then binary messages, and finally keyword- 
based messages. 
\item  Rule 2. As in mathematics, expressions in parentheses are executed before any others. 
\item Rule 3. When messages are of the same type, the order of execution is from left to right. 
\item In Smalltalk, mathematical message selectors such as + and * have the same priority, 
and therefore * does not have priority over +. You should use parentheses to ensure that 
your mathematical expressions are executed in the proper order. 

\end{itemize}

\ifx\wholebook\relax\else
    \end{document}
\fi

%%% Local Variables:
%%% coding: utf-8
%%% mode: latex
%%% TeX-master: t
%%% TeX-PDF-mode: t
%%% ispell-local-dictionary: "english"
%%% End:
